\documentclass[dvipdfmx, uplatex]{jsarticle}
\usepackage[dvipdfmx]{graphicx}
\usepackage{amsmath}
\usepackage{amssymb}
\usepackage{amsfonts}
\usepackage{amsthm}
\usepackage{mathrsfs}
\usepackage{bm}
\usepackage{braket}
\usepackage{xr}
\externaldocument{exercise_lattice_to_extra_dimension}

\usepackage{silence}
\WarningFilter*{hyperref}{Token not allowed in a PDF string}
\usepackage[dvipdfmx]{hyperref}
\usepackage{pxjahyper}
\allowdisplaybreaks[4]

\usepackage[left=25truemm,top=20truemm,bottom=20truemm,right=25truemm]{geometry}

\newcommand{\bk}{\bm{k}}
\newcommand{\bx}{\bm{x}}

\theoremstyle{definition}
\newtheorem*{sol}{解答}

\title{解答：格子ランダムウォークから余剰次元上の拡散へ\\
\large ―― \texttt{exercise\_lattice\_to\_extra\_dimension.tex} の解答 ――}
\author{}
\date{}

\begin{document}
\maketitle
\vspace{-6zh}

記号・式番号・問題番号は問題ファイル \texttt{exercise\_lattice\_to\_extra\_dimension.tex} のものを
踏襲する（第 I 部 (1)--(9)，第 II 部 (10)--(17)）。\textbf{各問の解（採点用）}は太字（\boxed）で示す。
以下，1 次元熱核を $G(x,t):=\dfrac{1}{\sqrt{4\pi Dt}}\,e^{-x^2/4Dt}$ と書く。

%======================================================================
\section{第 I 部の解答（格子ランダムウォークから拡散方程式へ）}
%======================================================================

\begin{sol}
\textbf{(1) 初期条件。}原点に局在ゆえ $\boxed{\,P_0(i,j)=\delta_{i,0}\,\delta_{j,0}\,}$
（時刻 $0$ で確率 $1$ が原点 $(0,0)$ に集中し，他の格子点では $0$）。

\par\smallskip
\textbf{(2) master 方程式。}時刻 $n+1$ に $(i,j)$ にいるのは，留まった（確率 $1-p$）か，隣接 4 点から
移ってきた（各 $\tfrac p4$）場合だから
\[
    \boxed{\,P_{n+1}(i,j)=(1-p)P_n(i,j)+\frac{p}{4}\bigl[P_n(i{+}1,j)+P_n(i{-}1,j)+P_n(i,j{+}1)+P_n(i,j{-}1)\bigr]\,}.
\]
全格子点で和をとると右辺 $=(1-p)\sum P_n+\tfrac p4\cdot4\sum P_n=\sum P_n$ ゆえ $\sum_{i,j}P_n$ は不変，
$\sum_{i,j}P_0=1$ より $\boxed{\sum_{i,j}P_n(i,j)=1}$（確率保存）。

\par\smallskip
\textbf{(3) 特性関数。}(2) に $\sum_{i,j}e^{i\bk\cdot\bx}$ を掛けて和をとり，添字をずらすと
（例：$\sum e^{i\bk\cdot\bx}P_n(i{+}1,j)=e^{-ik_xa}\hat P_n$）
$\hat P_{n+1}=\bigl[(1-p)+\tfrac p4(e^{ik_xa}+e^{-ik_xa}+e^{ik_ya}+e^{-ik_ya})\bigr]\hat P_n$，すなわち
\[
    \boxed{\,\lambda(\bk)=1-p+\frac{p}{2}\bigl(\cos k_xa+\cos k_ya\bigr)\,}.
\]
$\hat P_0(\bk)=\sum\delta_{i,0}\delta_{j,0}e^{i\bk\cdot\bx}=1$ ゆえ $\boxed{\,\hat P_n(\bk)=\lambda(\bk)^n\,}$。

\par\smallskip
\textbf{(4) 平均二乗変位。}1 ステップで $\Delta x=\pm a$ が各確率 $\tfrac p4$，それ以外は $\Delta x=0$ ゆえ
$\langle\Delta x^2\rangle=\boxed{\dfrac{pa^2}{2}}$（$\langle\Delta y^2\rangle$ も同じ）。各ステップ独立・平均 $0$ で
分散は加法的だから
\[
    \boxed{\,\langle x^2+y^2\rangle_n=n\,p\,a^2\,}.
\]

\par\smallskip
\textbf{(5) 素朴な極限は自明。}$t=n\tau$ 固定（$n=t/\tau$）ゆえ物理的広がりは
$\langle x^2+y^2\rangle=npa^2=\dfrac{t}{\tau}\,p\,a^2=p\,t\,\dfrac{a^2}{\tau}$。
\begin{itemize}
\item[(i)] $\tau$ 固定で $a\to0$：$n=t/\tau$ 固定，$\langle r^2\rangle=npa^2\to\boxed{0}$。有限歩数で歩幅 $0$ ゆえ
分布は原点に潰れる（$\delta$ のまま，拡散しない）。自明。
\item[(ii)] $a$ 固定で $\tau\to0$：$n=t/\tau\to\infty$，$\langle r^2\rangle\to\boxed{\infty}$。無限歩数で広がりが発散し，
規格化された極限分布を持たない（どこでも $0$）。自明。
\end{itemize}
一方だけの極限ではダメで，両者を結びつける必要がある。

\par\smallskip
\textbf{(6) diffusive scaling と $D$。}$\langle r^2\rangle=p\,t\,\dfrac{a^2}{\tau}$ を有限非ゼロに保つには，
$a,\tau\to0$ で $\boxed{\dfrac{a^2}{\tau}=\text{一定}}$（すなわち $\tau\propto a^2$，$a\sim\sqrt\tau$）に保てばよい
（diffusive scaling）。このとき $\langle x^2+y^2\rangle=4Dt$ とおいて拡散係数を定義すると
$4D=p\,\dfrac{a^2}{\tau}$ ゆえ
\[
    \boxed{\,D=\frac{p\,a^2}{4\tau}\,}.
\]

\par\smallskip
\textbf{(7) 連続極限の方程式と解。}

\textit{(a)} (2) を差の形にすると
$P_{n+1}(i,j)-P_n(i,j)=\dfrac p4\bigl[(P_{i+1}-2P_i+P_{i-1})+(P_{j+1}-2P_j+P_{j-1})\bigr]$。
右辺の 2 階差分は $P(x{\pm}a)=P\pm a\partial_xP+\tfrac{a^2}2\partial_x^2P\pm\cdots$ より
$P_{i+1}-2P_i+P_{i-1}=a^2\partial_x^2P+O(a^4)$（奇数階は相殺）。左辺は
$P(t{+}\tau)-P(t)=\tau\partial_tP+O(\tau^2)$。よって
\[
    \tau\,\partial_t P=\frac{pa^2}{4}\nabla^2P+O(a^4,\tau^2)
    \ \xrightarrow[\ a^2/\tau\ \text{一定}\ ]{}\
    \boxed{\,\frac{\partial P}{\partial t}=D\,\nabla^2 P\,}\quad\Bigl(D=\tfrac{pa^2}{4\tau}\Bigr).
\]

\textit{(b)} $\cos k_xa=1-\tfrac12(k_xa)^2+\cdots$ より $\lambda(\bk)=1-\tfrac{pa^2}{4}|\bk|^2+\cdots$。
$n=t/\tau$ ゆえ
$\hat P_n=\lambda^n=\exp[n\ln\lambda]\to\exp\!\bigl[-\tfrac{t}{\tau}\tfrac{pa^2}{4}|\bk|^2\bigr]
=\boxed{\,e^{-D t|\bk|^2}\,}$。逆 Fourier 変換（各次元 Gauss 積分）して
\[
    P(x,y,t)=\int\frac{d^2k}{(2\pi)^2}e^{-i\bk\cdot\bx}e^{-Dt|\bk|^2}
    =\boxed{\,\frac{1}{4\pi Dt}\exp\!\Bigl(-\frac{x^2+y^2}{4Dt}\Bigr)\,}\quad(\text{2 次元 Gauss 分布}).
\]

\textit{(c)} この $P$ は $\partial_tP=D\nabla^2P$ を満たす（$\hat P=e^{-Dt|\bk|^2}$ が $\partial_t\hat P=-D|\bk|^2\hat P$ を
満たすことから）。Gauss 分布より $\langle x^2\rangle=\langle y^2\rangle=2Dt$ ゆえ
$\boxed{\langle x^2+y^2\rangle=4Dt}$。これは (4) の $npa^2=\tfrac t\tau pa^2$ に (6) の $D=\tfrac{pa^2}{4\tau}$ を入れた
$4Dt$ に一致する。

\par\smallskip
\textbf{(8) マンハッタン距離と等確率面。}

\textit{(a)} $(0,0)\to(i,j)$ の最小ステップ数は軸方向にしか動けないので $\boxed{|i|+|j|}$
（マンハッタン距離 $=L^1$）。もし最短経路だけが効くなら $P\sim e^{-\beta(|x|+|y|)/a}$ 型となり，
等確率面は $\boxed{|x|+|y|=\text{一定}}$ ＝\emph{菱形}（$L^1$ 球）になる。

\textit{(b)} 終点 $x=(2k-N)a$ の経路数は $\Omega=\dbinom{N}{k}$。$k=N/2+\delta$（$\delta=x/2a$）とおき
Stirling より
\begin{align*}
    \ln\binom{N}{k}&\simeq N\ln2-\frac{2\delta^2}{N}-\frac12\ln\frac{\pi N}{2}\\
    &=N\ln2-\frac{x^2}{2Na^2}-\frac12\ln\frac{\pi N}{2}.
\end{align*}
よってエントロピー欠損 $\boxed{\Delta S=\ln\Omega_{\max}-\ln\Omega=\dfrac{x^2}{2Na^2}}$（$x$ の 2 次形式）。
$\Omega\propto e^{-x^2/2Na^2}$ ゆえ分散 $\boxed{\sigma_x^2=Na^2}$（$=\langle\Delta x^2\rangle\cdot N=a^2\cdot N$，
1 ステップ $\pm a$）。

\textit{(c)} (7)(b) で連続極限の特性関数は $\hat P_n\to e^{-Dt k_x^2}e^{-Dt k_y^2}$ と\emph{積に分解}した。
ゆえに $P(x,y,t)=G(x,t)\,G(y,t)\propto e^{-(x^2+y^2)/4Dt}$ であり，
\[
    P=\text{一定}\iff\boxed{x^2+y^2=\text{一定}}\quad(\text{円，}L^2\text{球}).
\]
分散は $\boxed{\sigma_x^2=\sigma_y^2=2Dt}$（等方）。

\textit{(d)} 確率は\emph{最短経路の長さ}（``エネルギー''）ではなく\emph{経路数＝多重度}
$\Omega=e^{S}$（``エントロピー''）で決まる：無偏りの歩みでは長さの等しい経路の重みに差はなく，
終点の確率はそこへ至る経路数に比例する。(b) のとおり Stirling によりその対数（エントロピー）は
終点座標の 2 次形式になり，しかも $x,y$ 方向の 1 ステップ分散が等しく
（$\langle\Delta x^2\rangle=\langle\Delta y^2\rangle$）相関がない（$\langle\Delta x\Delta y\rangle=0$）ため，
2 次形式は等方的な $x^2+y^2$ となる（中心極限定理）。これが等確率面が円（$L^2$）であって
菱形（$L^1$）でない理由である。菱形は，長い経路を指数的に罰する``エネルギー''項があって最短経路が
支配する場合（低温・大偏差領域）にのみ現れる。無偏り拡散ではエントロピーが支配し，格子の離散対称性
$C_4$ は連続極限で\textbf{回転対称性 $O(2)$ に創発的に拡大}する。この等方的な異方的スケーリング
（$x\to bx,\ t\to b^2t$，動的指数 $z=2$）の構造は，第 II 部のスケーリング次元の議論に引き継がれる。

\par\smallskip
\textbf{(9) 考察：微分の階数。}

\textit{(a)} 各向きへの移動確率が左右・上下対称ゆえ 1 ステップの平均変位は $\langle\Delta\bx\rangle=\bm0$。
このため 1 階差分（$\to\partial_x$，移流項）は相殺し，最低次で残るのは 2 階差分（$\to\partial_x^2$）である。
言い換えれば，消えない最低次の変位モーメントが $\langle\Delta x^2\rangle\sim a^2$（2 次）だから，空間微分は 2 階になる。

\textit{(b)} 連続極限の主要バランスは（時間 1 階）$\times\tau\ \sim\ $（空間 2 階）$\times a^2$，すなわち
$\tau\partial_tP\sim a^2\nabla^2P$。これが有限であるためには $\tau\sim a^2$（時間 $\sim$ 長さ$^2$），まさに
diffusive scaling である。つまり「時間 1 階・空間 2 階」という階数は $\tau\propto a^2$ と表裏一体。
時間が 1 階なのは，より高次の $\tau^2\partial_t^2\sim a^4$ がこのスケーリングで相対的に消えるからである。

\textit{(c)} 偏りがあると $\langle\Delta x\rangle\sim a\neq0$ となり，1 階差分が残って\emph{1 階}の空間微分項
（移流項）$\sim a\,\partial_xP$ が現れる。バランスは $\tau\partial_t\sim a\,\partial_x$ すなわち
$\boxed{\tau\propto a}$（ballistic／移流的，時間 $\sim$ 長さ）。このとき主要方程式は\emph{1 階}の移流方程式
$\boxed{\partial_tP=-v\,\partial_xP}$（$v=\lim\frac{\langle\Delta x\rangle}{\tau}$）となり拡散項
（$\sim a^2\nabla^2/\tau\sim a\to0$）は subleading。拡散を見るには移動座標系に移り，残差に diffusive scaling を
施す必要がある。\emph{消えない最低次の変位モーメントの次数が，連続極限の空間微分の階数と $\tau$ のスケール冪を決める}
――1 次（偏り）なら $\tau\sim a$ で移流，2 次（無偏り）なら $\tau\sim a^2$ で拡散，という対応である。
\end{sol}

%======================================================================
\section{第 II 部の解答（余剰次元上の拡散と次元簡約）}
%======================================================================

\begin{sol}
\textbf{(10) Fourier モード分解。}周期性 $y\sim y+2\pi R$ より $e^{ik_y\cdot2\pi R}=1$，すなわち
$\boxed{k_y=\dfrac{m}{R}}$（$m\in\mathbb{Z}$）。$\partial_y^2\to-(m/R)^2$ を 2 次元方程式に代入すると，各モードは
\[
    \boxed{\,\frac{\partial P_m}{\partial t}=D\,\frac{\partial^2 P_m}{\partial x^2}-\frac{Dm^2}{R^2}\,P_m\,}
\]
を満たす。$x$ 方向の拡散定数は $\boxed{D}$（不変），$m$ 依存の減衰率は $\boxed{\Gamma_m=\dfrac{Dm^2}{R^2}}$。
（``質量''$\mu_m=m/R$ で $\Gamma_m=D\mu_m^2$。）

\par\smallskip
\textbf{(11) ゼロモード。}$P_0(x,t)=\int_0^{2\pi R}P(x,y,t)e^{-i\cdot0\cdot y/R}dy=\int_0^{2\pi R}P\,dy$ は
$y$ で積分した周辺分布。(10) で $m=0$ ゆえ $\boxed{\partial_tP_0=D\partial_x^2P_0}$（純 1 次元拡散）。
点源 $P_0(x,0)=\int\delta(x)\delta(y)dy=\delta(x)$ より
\[
    \boxed{\,P_0(x,t)=G(x,t)=\frac{1}{\sqrt{4\pi Dt}}\,e^{-x^2/4Dt}\,}.
\]

\par\smallskip
\textbf{(12) KK 塔と寿命。}$P_m(x,0)=\int\delta(x)\delta(y)e^{-imy/R}dy=\delta(x)$ ゆえ，(10) の解は
\[
    \boxed{\,P_m(x,t)=e^{-\Gamma_m t}\,G(x,t)\,}.
\]
時定数 $\boxed{\tau_m=1/\Gamma_m=\dfrac{R^2}{Dm^2}}$。$|m|$ が小さいほど長寿命（$m=\pm1$ が最長）。
最低非ゼロモード $m=\pm1$ の寿命からクロスオーバー時刻
$\boxed{\,t^\ast:=\tau_1=\dfrac{R^2}{D}\,}$。

\par\smallskip
\textbf{(13) 全解と 2 表示。}$P=\sum_m\frac{1}{2\pi R}P_m e^{imy/R}=G(x,t)\,K(y,t)$，ただし円上の熱核 $K$ は
\begin{align*}
    \text{(i) モード和：}&\ \boxed{K(y,t)=\frac{1}{2\pi R}\sum_{m\in\mathbb{Z}}e^{imy/R}\,e^{-Dm^2t/R^2}},\\
    \text{(ii) 像和：}&\ \boxed{K(y,t)=\frac{1}{\sqrt{4\pi Dt}}\sum_{w\in\mathbb{Z}}e^{-(y-2\pi Rw)^2/4Dt}}.
\end{align*}
両者は Poisson 和で等しい。指数は (i) が $e^{-Dm^2t/R^2}$，(ii) が $e^{-(2\pi Rw)^2/4Dt}$ ゆえ，
\textbf{大 $t$ ではモード和 (i)}（$m=0$ 以外急減衰，数項で十分），\textbf{小 $t$ では像和 (ii)}
（$w=0$ 以外急減衰）が速く収束する。

\par\smallskip
\textbf{(14) 低エネルギー（$t\gg t^\ast$）。}(13)(i) で $m\neq0$ は $e^{-Dm^2t/R^2}=e^{-(m^2)t/t^\ast}\to0$。
生き残るのは $\boxed{m=0}$ のみで
\[
    \boxed{\,P(x,y,t)\simeq\frac{1}{2\pi R}\,G(x,t)=\frac{1}{2\pi R}\frac{1}{\sqrt{4\pi Dt}}e^{-x^2/4Dt}\,}
\]
――$y$ に一様，$x$ 方向の純 1 次元拡散（余剰次元は見えない）。最初の補正は $m=\pm1$ の項
$\frac{2}{2\pi R}e^{-Dt/R^2}\cos(y/R)\,G(x,t)$ であり，その大きさは $\boxed{\propto e^{-t/t^\ast}}$
（指数的に抑制）。

\par\smallskip
\textbf{(15) 高エネルギー（$t\ll t^\ast$）。}$y$ 方向の広がり $\sqrt{2Dt}\ll\sqrt{2}\,R<2\pi R$ ゆえ
粒子は円を一周しておらず，(13)(ii) で $\boxed{w=0}$ の項のみ支配的：
$K\simeq\frac{1}{\sqrt{4\pi Dt}}e^{-y^2/4Dt}$。よって
\[
    \boxed{\,P(x,y,t)\simeq\frac{1}{4\pi Dt}\,e^{-(x^2+y^2)/4Dt}\,}
\]
――自由 2 次元 Gauss（コンパクト性が見えない）。(13)(i) で実効的に効くのは $Dm^2t/R^2\lesssim1$，
すなわち $|m|\lesssim R/\sqrt{Dt}$ の項で，個数は $\boxed{N\sim\dfrac{R}{\sqrt{Dt}}\gg1}$（多数のモード）。

\par\smallskip
\textbf{(16) スケーリング次元と有効次元。}

\textit{(a)} $\partial_t\to b^{-z}\partial_t$，$\nabla^2\to b^{-2}\nabla^2$ ゆえ両辺が同じ冪になるのは
$\boxed{z=2}$（動的指数 $2$）。このとき $D$ は\fbox{不変（無次元・marginal）}。

\textit{(b)} $\int P\,d^dx=1$ かつ $x\to bx$（$d^dx\to b^d d^dx$）ゆえ $\boxed{P\to b^{-d}P}$。
自己相似形 $P=t^{-\alpha}\Phi(x/\sqrt{4Dt})$ が $t\to b^2t,\ x\to bx$ で不変なら，
プレファクター $t^{-\alpha}\to b^{-2\alpha}t^{-\alpha}$ が $b^{-d}$ と一致，すなわち $\boxed{\alpha=d/2}$。
$d=1$ で $\alpha=\tfrac12$（$t^{-1/2}$），$d=2$ で $\alpha=1$（$t^{-1}$）。

\textit{(c)(d)} $P(0,0,t)=G(0,t)\,K(0,t)$，$G(0,t)=\dfrac{1}{\sqrt{4\pi Dt}}\propto t^{-1/2}$。
\begin{itemize}
\item[$\bullet$] \textbf{長時間} $t\gg t^\ast$：$\sum_m e^{-Dm^2t/R^2}$ は $\boxed{m=0}$ の項のみ残り
$K(0,t)\to\dfrac{1}{2\pi R}$（定数）。ゆえ $P(0,0,t)\propto\boxed{t^{-1/2}}$，$\boxed{d_{\mathrm{eff}}=1}$。
効くモードは 1 個（最低 KK モード）。
\item[$\bullet$] \textbf{短時間} $t\ll t^\ast$：$\displaystyle\sum_m e^{-Dm^2t/R^2}\simeq\int dm\,e^{-Dm^2t/R^2}
=R\sqrt{\pi/Dt}$ ゆえ $K(0,t)\simeq\dfrac{1}{2\pi R}R\sqrt{\pi/Dt}=\dfrac{1}{\sqrt{4\pi Dt}}\propto t^{-1/2}$。
よって $P(0,0,t)\propto\boxed{t^{-1}}$，$\boxed{d_{\mathrm{eff}}=2}$。効くモードは $N\sim R/\sqrt{Dt}\gg1$ 個。
\end{itemize}
有効次元 $d_{\mathrm{eff}}=-2\times(\text{$P(0,0,t)$ の $t$ 冪})$ は，効く KK モード数が多い UV（短時間）で
$2$，最低モードだけが残る IR（長時間）で $1$ となる――\textbf{有効次元が $2\to1$ へ流れる}
（次元簡約のスケーリング的表現）。第 I 部 (8) の $C_4\to O(2)$ 創発（$z=2$）と合わせ，連続極限では
対称性が増え・有効次元が走る，という二つの``見かけ''の変化が現れている。

\par\smallskip
\textbf{(17) 考察。}

\textit{(a)} 指数 $e^{-Dm^2t/R^2}$（モード和）と $e^{-(2\pi Rw)^2/4Dt}$（像和）は，無次元量 $Dt/R^2$ について
互いに逆。$t\gg t^\ast$（$Dt/R^2\gg1$）ではモード和の非ゼロ項が消え 1 項（$m=0$）で済む＝(14)。
$t\ll t^\ast$（$Dt/R^2\ll1$）では像和の非ゼロ項が消え 1 項（$w=0$）で済む＝(15)。同じ関数の
\emph{速く収束する側}が極限ごとに入れ替わる（Poisson 和の双対性）。

\textit{(b)} 余剰次元（半径 $R$ の円）は $\Gamma_m=Dm^2/R^2$ を持つモードの塔を生む。$\Gamma_m$ は KK 質量
$\mu_m=m/R$ の 2 乗に拡散係数を掛けたもの（$\Gamma_m=D\mu_m^2$）で，KK 質量 $m/R$ に対応する。
ここで\textbf{逆時間 $1/t$ がエネルギースケール}，\textbf{$1/t^\ast=D/R^2$（最低 KK ``質量'' のスケール）が
コンパクト化スケール}に対応する。$1/t\ll1/t^\ast$（低エネルギー）では質量ゼロのゼロモードのみが効き
有効理論は 1 次元；$1/t\gg1/t^\ast$（高エネルギー）では塔全体が効き 2 次元が回復する
（(16) で見たスケーリング次元の流れ $d_{\mathrm{eff}}:2\to1$ はこの言い換え）。

\textit{(c)} (13) の双対は theta 関数のモジュラー変換（Poisson 和）であり，モジュラー的パラメータは
無次元の $Dt/R^2$。運動量 $m$（間隔 $1/R$）と巻きつき $w$（間隔 $2\pi R$）を入れ替える双対では，
次元解析上 $R$ は\textbf{$\tilde R\sim Dt/R$}に置き換わる（$R\,\tilde R\sim Dt$，すなわち拡散長の 2 乗 $Dt$ が
弦理論の $\alpha'$ に相当する不変スケール）。自己双対点は $R^2\sim Dt$，すなわち $t\sim t^\ast$ である。
\end{sol}

\end{document}
